% =====================================================================================
% book/preamble.tex — shared by book.tex and every standalone chapter driver.
%
% Compiled with xelatex (Unicode straight through: €, τ, ≈ appear in the source and in
% the notebook-generated captions, and must typeset, not vanish into a blank box).
% =====================================================================================

% ---------------------------------------------------------------- fonts
% The TeX Gyre faces ship with TeX Live but are not installed system-wide on macOS, so they
% are loaded by FILE NAME through kpathsea rather than by family name through fontconfig.
% (Loading by family name — \setmainfont{TeX Gyre Pagella} — fails here with "cannot be found".)
\usepackage{fontspec}
\setmainfont{texgyrepagella}[
  Extension = .otf, UprightFont = *-regular, BoldFont = *-bold,
  ItalicFont = *-italic, BoldItalicFont = *-bolditalic]
\setsansfont{texgyreheros}[
  Extension = .otf, UprightFont = *-regular, BoldFont = *-bold,
  ItalicFont = *-italic, BoldItalicFont = *-bolditalic, Scale = MatchLowercase]
\setmonofont{texgyrecursor}[
  Extension = .otf, UprightFont = *-regular, BoldFont = *-bold,
  ItalicFont = *-italic, BoldItalicFont = *-bolditalic, Scale = MatchLowercase]
\usepackage{unicode-math}
\setmathfont{texgyrepagella-math.otf}
\usepackage{microtype}

% ---------------------------------------------------------------- maths
\usepackage{amsmath}
\usepackage{amsthm}
\numberwithin{equation}{chapter}   % (9.3), not (137) — a book numbers by chapter
\DeclareMathOperator*{\argmin}{arg\,min}
\DeclareMathOperator*{\argmax}{arg\,max}

% Named, numbered assumptions. Identification lives or dies on these, so they are
% first-class objects in this book, not bullets: they get numbers, and the text refers
% to them by number ("Assumption 9.2 is the one the placebos cannot test").
\theoremstyle{definition}
\newtheorem{assumption}{Assumption}[chapter]
\newtheorem{proposition}{Proposition}[chapter]
\theoremstyle{remark}
\newtheorem*{remark}{Remark}

% ---------------------------------------------------------------- floats, tables, figures
\usepackage{graphicx}
\usepackage{booktabs}
\usepackage{longtable}
% max width=\linewidth SHRINKS a table that would overrun the text block and leaves a narrow one
% alone (unlike \resizebox{\linewidth}, which would also blow small tables up). Applied by
% build.py to every generated table, so a wide results table can never run into the margin.
\usepackage{adjustbox}
\usepackage{caption}
\captionsetup{font=small, labelfont=bf, width=.92\linewidth}
\usepackage{float}
\usepackage[table]{xcolor}
\usepackage{enumitem}

% ---------------------------------------------------------------- house colours
\definecolor{SchemeDark}{HTML}{1F4E66}
\definecolor{SchemeBlue}{HTML}{2C7FB8}   % the model / our estimate  (matches cmp.plots)
\definecolor{SchemeOrange}{HTML}{D95F0E} % the naive baseline / the cost line
\definecolor{SchemeGrey}{HTML}{999999}

% ---------------------------------------------------------------- links, refs, bibliography
\usepackage[colorlinks=true, linkcolor=SchemeDark, citecolor=SchemeDark,
            urlcolor=SchemeDark, breaklinks=true]{hyperref}
\usepackage[nameinlink]{cleveref}
\crefname{assumption}{Assumption}{Assumptions}
\Crefname{assumption}{Assumption}{Assumptions}
\usepackage[backend=biber, style=authoryear, natbib=true, maxcitenames=2,
            uniquelist=false, giveninits=true, doi=false, isbn=false, url=false]{biblatex}

% ---------------------------------------------------------------- money and per cent
% Numbers themselves are NEVER written here: they arrive as macros from build/macros.tex,
% which is generated from the executed notebooks. These commands only carry the typography.
\newcommand{\keur}[1]{\mbox{\texteuro#1\,k}}   % \keur{\nbSevenClTotal}  ->  €318 k
\newcommand{\eur}[1]{\mbox{\texteuro#1}}
\newcommand{\kusd}[1]{\mbox{\$#1\,k}}
\newcommand{\usd}[1]{\mbox{\$#1}}
\newcommand{\pc}[1]{\mbox{#1\,\%}}
\newcommand{\pcof}[1]{\mbox{#1\,\%}}

% ---------------------------------------------------------------- notation
\newcommand{\Prob}{\operatorname{\mathbb{P}}}
\newcommand{\E}{\operatorname{\mathbb{E}}}
\newcommand{\Var}{\operatorname{Var}}
\newcommand{\sd}{\operatorname{sd}}
\newcommand{\simplex}{\Delta^{J-1}}
\newcommand{\Yobs}{Y}
\newcommand{\code}[1]{\texttt{#1}}

% ---------------------------------------------------------------- the notebook pointer box
% Every chapter closes with one of these: where the runnable code is, which seed, FAST vs FULL.
% Deliberately a small note, not a section — the book is the artifact, the notebook is the lab.
\newcommand{\notebooknote}[1]{%
  % ragged right, not \sloppy: the note is dense with long unbreakable \code{} paths, and
  % justifying around them opens word gaps you can drive a bus through.
  \par\medskip\noindent\begingroup\small\raggedright
  \textcolor{SchemeGrey}{\rule{\linewidth}{0.4pt}}\par\smallskip
  \textbf{The runnable notebook.} #1\par
  \endgroup\medskip}

% ---------------------------------------------------------------- generated inputs
% build/macros.tex   : one \newcommand per number the notebooks emitted (cmp.report.value)
% build/tables/*.tex : booktabs floats                                  (cmp.report.table)
% build/floats/*.tex : figure floats, captions and all                  (cmp.report.figure)
%
% If a macro is missing, xelatex raises "Undefined control sequence" and book/build.py
% refuses to ship the PDF. That is the whole point: a hole is louder than a stale number.
\graphicspath{{build/}{build/figures/}{../build/}{../build/figures/}}
